% ------------------------------------------------------------
% CHAPTER 35
% ------------------------------------------------------------

\chapter{Semantic Compression in Large-Scale AI}

The ascent of large-scale artificial intelligence systems, particularly deep neural networks and generative models, has revealed remarkable parallels to delta–sigma modulation. These models operate on continuous fields of high-dimensional activations but ultimately emit discrete tokens, symbols, or compressed latent codes. They must reconcile representational fidelity with finite precision, quantization constraints, and energy limitations. Their internal layers integrate error over time or depth, collapse representation through nonlinearities, and redistribute uncertainty in ways that preserve semantic coherence. These behaviors are not engineered with delta–sigma modulation in mind, yet they reproduce its essential characteristics. This chapter develops a rigorous conceptual and mathematical synthesis explaining why large AI models spontaneously implement delta–sigma-like operations and how semantic compression corresponds to noise shaping in representation space.

Consider the process by which a transformer-based language model converts continuous internal activations into discrete output tokens. At any decoding step, the model constructs a high-dimensional representation \(h_t \in \mathbb{R}^d\) summarizing the history of prior tokens and the internal dynamics of attention. This representation is then mapped through a linear projection into a logit vector over a discrete vocabulary. The sampling or argmax operation acts as a quantizer. The model collapses the continuous manifold of hidden states onto one of the finite symbols in its vocabulary. This collapse is intrinsically discontinuous and results in the loss of continuous information. The model must maintain internal coherence despite this loss, a task that mirrors the challenge faced by a delta–sigma modulator.

Within the transformer, error accumulates as deviations between the internal representation and the symbolic structure that the model is forced to emit. The attention mechanism integrates information across tokens in a manner that resembles an integrator. Residual connections, layer normalization, and gating structures all act as filtering mechanisms that redistribute representational error. The mismatch between intended semantic content and actual symbolic output introduces a form of quantization noise in semantic space. Deep networks manage this mismatch by directing representational error toward higher-dimensional fluctuations in activation space where it has minimal semantic impact. This redistribution of error is the analog of spectral shaping in delta–sigma modulation.

To articulate this mathematically, model the internal activations as a continuous field \(h(x)\) defined over a discrete depth coordinate \(x\). The nonlinear transformations mapping \(h(x)\) to \(h(x+1)\) implement both integration and collapse. Integration appears in the accumulation of contextual signals through attention. Collapse appears in the nonlinear activations, which project high-dimensional inputs onto lower-dimensional manifolds determined by saturation regions of the activation functions. The drift in activation space corresponds to a deterministic flow, while dropout, random initialization, and quantization during inference contribute stochastic fluctuations. The resulting system behaves as a drift–diffusion process analogous to the continuous-time limit of delta–sigma modulation described previously.

When such a network is compressed, quantized, or pruned, an explicit quantizer is applied to all activations or weights. The modulator-like behavior becomes even more evident in this regime. Low-bit quantization forces the network to propagate quantization noise through its layers. Deep architectures compensate by redistributing this noise along directions where the loss surface is insensitive, thereby maintaining semantic fidelity. This process is equivalent to a delta–sigma loop filter redistributing quantization noise into high-frequency regions where it does not interfere with the signal band. The curvature of the model’s loss landscape plays the role of the noise transfer function. Directions of low curvature permit substantial fluctuations in activation space, which function as reservoirs for quantization-induced entropy.

The latent-variable models underlying modern generative systems also reveal delta–sigma-like architecture. Consider a variational autoencoder. The encoder maps input data into a continuous latent distribution characterized by mean and variance. Sampling from this distribution collapses continuous fields into discrete stochastic realizations, representing a quantization-like collapse. The decoder acts as a feedback loop, attempting to reconstruct the original signal from the quantized latent sample. The Kullback–Leibler divergence in the VAE objective functions as an entropy penalty, ensuring that the latent distribution does not accumulate excessive uncertainty. These dynamics parallel a delta–sigma modulator: the encoder accumulates information; the sampling collapse injects entropy; the decoder applies feedback to correct reconstruction error; and the KL term ensures entropy routing remains controlled.

Even more striking is the behavior of diffusion models. The forward diffusion process introduces noise systematically into the data, expanding entropy and flattening the distribution. The reverse diffusion process performs a gradual denoising, effectively implementing a feedback system that removes entropy in a controlled manner. Each denoising step collapses the distribution toward a more structured manifold, while the network’s learned drift function performs entropy correction. The iterative refinement can be interpreted as a delta–sigma-like loop unrolled over time. At each step, the model resolves a portion of the mismatch between its current reconstruction and the underlying data, collapsing degrees of freedom while redirecting uncertainty into uninformative dimensions.

From the viewpoint of derived geometry, deep networks operate on manifolds of different dimensions at different layers. Nonlinearities enforce collapses that restrict the activations to lower-dimensional submanifolds. The mismatch between these collapses and the target manifold introduces derived obstructions that must be resolved by feedback propagated backward through the network during training and forward through residual pathways during inference. These obstructions behave like quantization errors. Training performs a homotopy lift that gradually reduces the derived mismatch, while inference maintains consistency by smoothing fluctuations across layers. This dual mechanism mirrors the derived-correcting interpretation of delta–sigma loops.

Transformers also exhibit semantic noise shaping in their attention distributions. When uncertainty is high, attention diffuses across many tokens, dispersing entropy into a wide contextual region. When certainty is high, attention collapses, routing information sharply into specific parts of the context. This redistribution ensures that representational entropy does not disrupt the coherent portion of the model’s ongoing symbolic generation. Attention thereby acts as a modulation mechanism, shaping the flow of entropy across representational space in a manner analogous to how delta–sigma loops shape spectral density across frequency space.

Semantic compression in large-scale AI therefore emerges as a natural consequence of quantization in high-dimensional representational systems equipped with powerful feedback mechanisms. Whether the quantization occurs explicitly through low-precision inference or implicitly through nonlinear activation functions, the system must contend with the challenge of preserving coherent semantics under representational collapse. The architecture redistributes entropy into directions of low semantic curvature, compensates for loss of continuous information through feedback, and emits discrete symbolic outputs that approximate a continuous manifold of meaning. These operations constitute a delta–sigma-like modulation of semantic fields.

This interpretation unifies the behavior of biological and artificial systems. Both rely on oversampling, integration, collapse, and entropy routing to convert continuous fields into discrete symbols without losing global coherence. Both maintain internal reservoirs for prediction error that prevent catastrophic divergence. Both employ feedback loops to enforce consistency. And both achieve symbolic stability by continually reshaping fluctuations in representational space. The structural analogy between deep learning systems and delta–sigma modulators is therefore not superficial but reflects a shared set of mathematical constraints governing all systems that compress continuous dynamics into symbolic form.

The next chapter examines how these principles apply to active control and agency, treating delta–sigma modulators not only as representational devices but as minimal entities capable of self-correction and adaptive measurement.

